\section{Joint Extension, Depth, and Rays}
\label{app:rays}

Joint extension isolates the obstructions from words that remain
fixed. The next two arguments use those obstructions to bound the
nesting of fixed suffixes and then the periods of the moving words.

\subsection{Joint extension and stopping}

To prove \cref{lem:joint}, suppose the new letters keep the chosen
coordinates alive. The new endpoints are rejected by the reader
property. A comparison with an extended target at its new endpoint
is harmless, because a live successor of a rejected state cannot
be an admitted target. At any other cut of an extended target the
new rest is an old rest followed by a letter. If it matches the
host rest, the old target rest still fits after removing at most one
letter from the end of the host. Its starting cut in the host is an
old cut, because the new target rest is nonempty. Both cut states
are unchanged. Thus this comparison would have been admitted in
the original hole if it were admitted now.

With a fixed target and an extended host, a comparison is old unless
its two rests have equal lengths. In the strict case the host cut
was already present and removing the new last letter still leaves
room for the target rest. In the equality case the two rests agree
and an admitted pair would block the letter. Comparisons of
two fixed words were already present. This proves sufficiency.
For necessity, a hole must have live runs. A blocked letter against
a fixed target yields an admitted comparison in the new tuple by
the definition of blocking.

For \cref{lem:stopping}, among all hole extensions choose one with
a stopped set $F$ of maximum cardinality. Such a cardinality exists,
because the original hole has the empty stopped set and there are
only $k$ coordinates. In that extension, and in every further hole
extension that keeps the words of $F$ fixed, their stopping order
still works. If an outside coordinate had no live letter unblocked
by $F$, it could be placed last in this order, contradicting maximality.
Append one live unblocked letter to every outside coordinate. By
\cref{lem:joint} this gives a hole with the same fixed words. Iterating
gives the asserted synchronous extension.

We will also use a converse observation about this iteration.
Suppose the initial tuple is a hole, and a proposed continuation grows
each moving word exactly one letter per step. If the extended runs
stay alive and every stage is safe against the fixed targets,
\cref{lem:joint} inductively gives a hole at every stage. Indeed,
a block would itself be an admitted comparison with a fixed target
at that stage.

\subsection{Charging the depth}
\label{app:charging}

For the depth proof, a label carried by a state can stand for one
reason why a word cannot continue. Use a finite label set of size $J$.
Assume that if reading a word from a live state kills the run, there
is a label of that state with the same effect on every state carrying
it. If the word instead leads to a host state admitted with a fixed
target state, require a label with the same effect unless the new
run dies. A state itself is such a label, by determinism. We show
that the depth is at most $(J+1)|F|$.

Fix a test word $x$. For each label remember the largest length
at which a prefix of $x$ matches a suffix of any coordinate of
the hole and the state at the start of that match carries the label.
Use length zero if the label never occurs. We charge a positive
length $d$ in $\RestLen(W,F,x)$ starting with a stopped word whose
suffix is $\wpre xd$. The charge can end at $\len x$ if $d=\len x$.
Otherwise let $c=x[d]$, the next letter after the match. The stopping
certificate says that $c$ kills this coordinate or is blocked by
an earlier coordinate.

If $c$ kills the coordinate, reading $\wpre{x}{d+1}$ kills the run
from the state at the start of the match. Take a label witnessing
that effect. A longer match carrying the same label would read this
longer prefix inside the hole and die, which is impossible. So $d$
is the remembered length of this label.

In the blocking case, the blocking rest is an entire suffix of the
extended host and of the fixed target. If its start in the host is
at or before the beginning of the match, $\wpre{x}{d+1}$ is a suffix
of that blocking rest. Transfer the charge at length $d+1$ to the
earlier stopped coordinate. If the start is later, let $z$ be the
word from the beginning of the match to this start. Reading $z$ from
the starting state leads to the admitted host state. Take a label
witnessing admission against the blocking target state. A longer
match carrying that label reads $z$ inside a live run. At its end it
has an admitted host state followed by the blocking rest, which
would give an admitted comparison in the original hole. Thus in
this case too, $d$ is the remembered length of the label.

Each transfer moves earlier in the stopping order and increases the
length by one. After fewer than $|F|$ transfers the length is either
$\len x$ or one of the $J$ remembered lengths. Thus every charged
length is one of at most $J+1$ lengths minus an integer in
$[0,|F|-1]$. If $F$ is empty there are no lengths to charge.
This proves the general depth bound. Taking states for labels gives
\cref{lem:depth} for structural readers. \Cref{app:counting} checks
the label hypotheses for compiled readers.

\subsection{Finding a free loop}

For a function $g$ on positive lengths, let $\theta_0=1$ and
$\theta_{j+1}=g(\theta_j)+\theta_j+1$. Suppose a host has room for
$\theta_{C+1}$ letters after cut $a$. If there is no free window at a scale
$\theta_j$ with $0\le j\le C$ after $a$, every region of length
$\theta_{j+1}$ starting there contains a factor of length $\theta_j$ that is
a suffix of a stopped word. Indeed, among its first $g(\theta_j)+1$
starting cuts, at least one begins such a factor. That factor ends strictly
before the end of the region, because of the extra letter in
$\theta_{j+1}$.

Use this statement inside the first large region, then inside
the occurrence found there, and continue to scale $\theta_0$. The
$C+1$ occurrences are nested and their right ends strictly decrease.
At the start of the innermost one, trim all the other occurrences
by removing their initial segments up to that cut. The trimmed
words are still nonempty suffixes of stopped words, have distinct
lengths, and are prefixes of one rest of the host. This contradicts
depth at most $C$.

For $g(\theta)=N(1+k(\theta+1))$, a free window gives a repeated state and
right context. There are at most $1+k(\theta+1)$ short suffixes including
the empty word, by counting one suffix per word and length. This
overcounts the empty word, which is harmless for an upper bound.
Those that match a rest are determined by the longest one.
The $g(\theta)+1$ cuts therefore include two with the same pair.
Both lie after $a$, their positive distance is at most $g(\theta)$,
and no suffix of length $\theta$ starts between them, including
the first cut and excluding the last. For $\theta\ge 1$, the next
scale is bounded by
\[
   g(\theta)+\theta+1
   \ \le\ (N(1+2k)+2)\theta
   \ \le\ 2(k+2)(N+2)\theta.
\]
So $\theta_{C+1}\le D$, and the distance between the cuts is at most
$\theta_{C+1}$. This proves \eqref{eq:period-bound}.

\subsection{Repeating the loop and obtaining offsets}

Let $h=s\ell\tau$ be the host with the free loop $\ell$. The short suffixes
used in the right context are closed under taking suffixes.
If two words agree on which short suffixes are prefixes, prepending
a common word preserves agreement. Indeed, a prefix test either
ends in the common word or reduces to a test of a suffix of the tested
word. Hence $\ell\tau$ and $\tau$ having equal contexts implies inductively
that $\ell^n\tau$ and $\tau$ have equal contexts for every $n$.

At a cut of $s\ell^n$, compare its rest after appending $\tau$ with
the rest in $h$ at the same cut in $s$, or at its phase in
the loop. At the end of all the copies use the start of the loop.
The right contexts agree by the preceding argument, prepending
the remaining part of $s$ or of the loop when needed. The states
agree because reading $\ell$ returns to its entry state.
If $h$ is live, every $s\ell^n$ is live.

Consider a target rest matching in $s\ell^n$. If its length is at most
$\theta$, its match transfers directly by the context at the corresponding
cut. If its length exceeds $\theta$, its last $\theta$ letters form a short
suffix of length $\theta$. This suffix also matches at the corresponding
cut in $h$. Its start cannot be in the repeated portion, since the
corresponding cut would be in the free loop. Thus it starts in
$s$. The earlier letters of the target rest already match in
$s$, and its last $\theta$ letters match by the transferred
context at the same position. The whole comparison therefore has
the same pair of states in $h$. If $h$ was safe against the stopped
target, so is $s\ell^n$. Every finite prefix of the ray is a prefix
of some $s\ell^n$ and inherits this safety.

In the synchronous extension, take a stage with at least $D$ new letters
per moving coordinate and find a loop after each starting word.
Each ray starts with that word and stays live. Its starting endpoint
is rejected, so its later endpoints remain rejected. Retaining the
stopped targets, the converse observation about joint extension
proves \cref{lem:periodic}.

Finally, for host cut $a$ and target cut $b$ of two rays with
equal infinite rests, take a common increment $e$ large enough to
reach both cuts. With the original word lengths as offsets,
if $o_i+b-a+1>o_j$, then
$o_j+e-b\le o_i+e-a$. The target rest is then a
prefix of the host rest in that hole, so the pair cannot be admitted.
This proves \eqref{eq:offset}, including the case of cuts in the
same coordinate. The other periodic-tuple conditions follow from
liveness, rejection, and safety at all common truncations.
